The author set the research question, the constraints of the study (own domains only, providers whose terms allow
publication) and the requirement that the result be a usable tool, and approved each protocol
before it was frozen. An AI coding agent (Claude Code, model Claude Opus), working in the author's project under the
author's direction, proposed the protocols, wrote the library and the experiment code, ran the experiments and the
analyses, and drafted this manuscript; a sub-agent compiled the bibliography, each entry verified against a primary
source. The simulated payloads are generated by code from public OpenAPI descriptions, not written by a model. No
human annotation was performed beyond the author's review. The work added in version 2, Study 3 and the software
from version 0.4 on, was carried out in the same way by AI coding agents under the author's direction: Claude Code
and, from release 1.2.0, Codex, which developed releases 1.2.0 to 1.3.1 and parts of 1.4 under the coordination of
a Claude Code session; the author decided its scope, approved its protocols and decided every publication.
In Study~4 (version 3) Codex drafted protocol v3 and built Part~A's worlds, prompts and rubric; Claude Code built
Part~B's tasks, ran the campaigns and the analyses and drafted the section. The answers were rated by an AI judge
(\texttt{\SfourJudge}) called by a script, never by the agent that ran the campaigns, and on a sample by Codex
(GPT-6.2 Astra) in a separate session, the agent that had drafted protocol v3 and its rubric and worked on the
library. The author chose the judges, approved each protocol and the change of judge before any score, and
authorised every paid run; the other deviations, recorded by the agent, are listed in Section~\ref{sec:study4}.
In version 4 Claude Code drafted protocols v3d and v3e, implemented the text previews and the LangGraph adapter,
ran the campaigns and the analyses of Parts~D and~E and drafted them; a sub-agent with its own context recomputed
Part~E's results from the traces. The author approved each protocol and authorised every paid run.
The code and the manuscript went through repeated reviews by several AI agents, among them Claude Code (Claude Opus
5.5) and Codex (GPT-6.2 Astra), in separate sessions or in sub-agents with their own context; the problems they
classed as blocking were fixed, those in the code with a regression test, and the others were fixed, documented as
limits of the software or of this study, or left to the author's decision. The author is responsible for the
decision to publish and for the final text.
